
\documentclass{article}
\usepackage{colortbl}
\usepackage{makecell}
\usepackage{multirow}
\usepackage{supertabular}

\begin{document}

\newcounter{utterance}

\centering \large Interaction Transcript for game `function\_detective', experiment `numbers\_passive\_examples', episode 7 with mistral{-}large{-}3{-}azure.
\vspace{24pt}

{ \footnotesize  \setcounter{utterance}{1}
\setlength{\tabcolsep}{0pt}
\begin{supertabular}{c@{$\;$}|p{.15\linewidth}@{}p{.15\linewidth}p{.15\linewidth}p{.15\linewidth}p{.15\linewidth}p{.15\linewidth}}
    \# & \multicolumn{2}{c}{Player} && \multicolumn{2}{c}{Game Master} \\
    \hline

    \theutterance \stepcounter{utterance}  
    & & & \multicolumn{4}{p{0.6\linewidth}}{
        \cellcolor[rgb]{0.9,0.9,0.9}{
            \makecell[{{p{\linewidth}}}]{
                \texttt{\tiny{[P1$\langle$GM]}}
                \texttt{I have a hidden Python function for you to figure out.} \\
\\ 
\texttt{Your job is to infer what it does and then submit a correct Python implementation.} \\
\\ 
\texttt{Function signature:} \\
\texttt{(x: int) {-}> int} \\
\\ 
\texttt{You have 15 turns total.} \\
\\ 
\texttt{I control which valid input{-}output examples you see. You can request them one at a time.} \\
\\ 
\texttt{You must make your best SOLVE attempt no later than the second{-}to{-}last turn.} \\
\texttt{Do not wait until the very end.} \\
\\ 
\texttt{You can use these actions:} \\
\texttt{{-} JSON format described at the end of the prompt} \\
\texttt{{-} SOLVE: ```python ... ```} \\
\\ 
\texttt{Do not add import statements in your guess.} \\
\texttt{Anything you might need is already available, so just write the function body in hypothesis and return the answer.} \\
\\ 
\texttt{Preloaded examples:} \\
\texttt{1. ({-}1, 0)} \\
\\ 
\texttt{Scoring:} \\
\texttt{{-} correctness is required for any non{-}zero score} \\
\texttt{{-} in interactive modes, solving with fewer information requests is better} \\
\\ 
\texttt{Input formatting reminders:} \\
\texttt{{-} strings must be quoted} \\
\texttt{{-} booleans must be True/False} \\
\texttt{{-} use commas between multiple values} \\
\texttt{{-} do not write any extra action tags in the same message} \\
\\ 
\texttt{Format example:} \\
\texttt{You write:} \\
\texttt{\{'current\_mode': <probing/confirming/don't know; choose whichever best describes your current state in the task>, 'input\_rationale': <brief description of why you chose this hypothesis>, 'hypothesis': <either the exact string "unknown" if you have no guess yet OR your current best guess as runnable Python matching the given signature, with a one{-}line docstring stating what it does. Never write placeholder code, if you cannot commit to a guess, use "unknown". Match the given signature.} \\
\texttt{e.g.} \\
\texttt{def f(x):} \\
\texttt{	"""Returns the absolute value of x."""} \\
\texttt{	return x if x >= 0 else {-}x} \\
\\ 
\texttt{Use "unknown" if you have no guess.>, 'confidence\_score': <score between 1{-}100. 100 means highest confidence in your hypothesis>\}} \\
\\ 
\texttt{Output only one valid json per turn.} \\
            }
        }
    }
     \\ \\

    \theutterance \stepcounter{utterance}  
    & \multicolumn{4}{p{0.6\linewidth}}{
        \cellcolor[rgb]{0.561,0.737,0.561}{
            \makecell[{{p{\linewidth}}}]{
                \texttt{\tiny{[P1$\rangle$GM]}}
                \texttt{```json} \\
\texttt{\{} \\
\texttt{    "current\_mode": "probing",} \\
\texttt{    "input\_rationale": "The preloaded example ({-}1, 0) suggests that the function returns 0 for negative inputs. I will test a positive input to see if it returns the input itself or something else.",} \\
\texttt{    "hypothesis": "unknown",} \\
\texttt{    "confidence\_score": 10} \\
\texttt{\}} \\
\texttt{```} \\
            }
        }
    }
    & & \\ \\

    \theutterance \stepcounter{utterance}  
    & & & \multicolumn{4}{p{0.6\linewidth}}{
        \cellcolor[rgb]{0.9,0.9,0.9}{
            \makecell[{{p{\linewidth}}}]{
                \texttt{\tiny{[P1$\langle$GM]}}
                \texttt{OUTPUT: (1000, 2)} \\
            }
        }
    }
     \\ \\

    \theutterance \stepcounter{utterance}  
    & \multicolumn{4}{p{0.6\linewidth}}{
        \cellcolor[rgb]{0.561,0.737,0.561}{
            \makecell[{{p{\linewidth}}}]{
                \texttt{\tiny{[P1$\rangle$GM]}}
                \texttt{```json} \\
\texttt{\{} \\
\texttt{    "current\_mode": "probing",} \\
\texttt{    "input\_rationale": "The preloaded example ({-}1, 0) and the new example (1000, 2) suggest the function might be counting the number of digits in the absolute value of x, but 1000 has 4 digits, not 2. Alternatively, it could be counting the number of trailing zeros in the binary representation, but 1000 in binary is 1111101000, which has 3 trailing zeros. Another possibility is the number of times x can be divided by 10 before becoming odd, but 1000 divided by 10 three times is 1 (odd), so that doesn't match. I will test another input to narrow it down.",} \\
\texttt{    "hypothesis": "unknown",} \\
\texttt{    "confidence\_score": 20} \\
\texttt{\}} \\
\texttt{```} \\
            }
        }
    }
    & & \\ \\

    \theutterance \stepcounter{utterance}  
    & & & \multicolumn{4}{p{0.6\linewidth}}{
        \cellcolor[rgb]{0.9,0.9,0.9}{
            \makecell[{{p{\linewidth}}}]{
                \texttt{\tiny{[P1$\langle$GM]}}
                \texttt{OUTPUT: (1, 0)} \\
            }
        }
    }
     \\ \\

    \theutterance \stepcounter{utterance}  
    & \multicolumn{4}{p{0.6\linewidth}}{
        \cellcolor[rgb]{0.561,0.737,0.561}{
            \makecell[{{p{\linewidth}}}]{
                \texttt{\tiny{[P1$\rangle$GM]}}
                \texttt{```json} \\
\texttt{\{} \\
\texttt{    "current\_mode": "probing",} \\
\texttt{    "input\_rationale": "The new example (1, 0) suggests that the function returns 0 for x=1, which contradicts the digit{-}counting hypothesis. The pattern so far: ({-}1, 0), (1, 0), (1000, 2). This looks like the function might be counting the number of trailing zeros in the decimal representation of x. For x=1, there are 0 trailing zeros; for x=1000, there are 3 trailing zeros, but the output is 2, which doesn't match. Alternatively, it could be the number of times x can be divided by 10 before reaching a number less than 100. For x=1000: 1000/10=100 (1), 100/10=10 (2), 10/10=1 (3) → but output is 2. This doesn't fit. Another possibility is the number of times x can be divided by 10 before the result is less than 100, minus 1. For x=1000: 1000/10=100 (1), 100/10=10 (2) → output is 2. For x=1:} \\
            }
        }
    }
    & & \\ \\

    \theutterance \stepcounter{utterance}  
    & & & \multicolumn{4}{p{0.6\linewidth}}{
        \cellcolor[rgb]{0.9,0.9,0.9}{
            \makecell[{{p{\linewidth}}}]{
                \texttt{\tiny{[P1$\langle$GM]}}
                \texttt{Game Violation: ParseError: Response is not valid JSON.} \\
\texttt{Output ONLY:} \\
\texttt{One JSON with confidence\_score, current\_mode, input\_rationale, and hypothesis keys ...} \\
\texttt{OR} \\
\texttt{SOLVE: ```python ... ```} \\
            }
        }
    }
     \\ \\

    \theutterance \stepcounter{utterance}  
    & \multicolumn{4}{p{0.6\linewidth}}{
        \cellcolor[rgb]{0.561,0.737,0.561}{
            \makecell[{{p{\linewidth}}}]{
                \texttt{\tiny{[P1$\rangle$GM]}}
                \texttt{```json} \\
\texttt{\{} \\
\texttt{    "current\_mode": "probing",} \\
\texttt{    "input\_rationale": "The examples ({-}1, 0), (1, 0), and (1000, 2) suggest the function counts the number of trailing zeros in the decimal representation of x, but capped at 2. For x=1000 (3 trailing zeros), the output is 2, which fits this hypothesis. For x=100, I predict the output would be 2. I will test this hypothesis with x=100.",} \\
\texttt{    "hypothesis": "unknown",} \\
\texttt{    "confidence\_score": 50} \\
\texttt{\}} \\
\texttt{```} \\
            }
        }
    }
    & & \\ \\

    \theutterance \stepcounter{utterance}  
    & & & \multicolumn{4}{p{0.6\linewidth}}{
        \cellcolor[rgb]{0.9,0.9,0.9}{
            \makecell[{{p{\linewidth}}}]{
                \texttt{\tiny{[P1$\langle$GM]}}
                \texttt{OUTPUT: ({-}2, 2)} \\
            }
        }
    }
     \\ \\

    \theutterance \stepcounter{utterance}  
    & \multicolumn{4}{p{0.6\linewidth}}{
        \cellcolor[rgb]{0.561,0.737,0.561}{
            \makecell[{{p{\linewidth}}}]{
                \texttt{\tiny{[P1$\rangle$GM]}}
                \texttt{```json} \\
\texttt{\{} \\
\texttt{    "current\_mode": "confirming",} \\
\texttt{    "input\_rationale": "The new example ({-}2, 2) breaks the trailing zeros hypothesis since {-}2 has no trailing zeros. However, the absolute value of {-}2 is 2, which has 0 trailing zeros, but the output is 2. This suggests the function might be related to the number of digits or some other property. Given the examples: ({-}1, 0), (1, 0), ({-}2, 2), (1000, 2), I now hypothesize the function returns the number of times the input can be divided by 2 before becoming odd (i.e., the exponent of 2 in the prime factorization of x). For {-}1: 0 (odd), for 1: 0 (odd), for {-}2: 1 (but output is 2, which doesn't fit). Alternatively, it could be the number of trailing zeros in the binary representation of x. For {-}1: binary is all 1s (0 trailing zeros), for 1: 1 (0 trailing zeros), for {-}2: binary is {-}10 (1 trailing zero, but output is 2). This doesn't fit. Another possibility: the number of times x can be divided by 10 before becoming less than 10 in absolute value, but for {-}2: {-}2/10 = {-}0.2} \\
            }
        }
    }
    & & \\ \\

    \theutterance \stepcounter{utterance}  
    & & & \multicolumn{4}{p{0.6\linewidth}}{
        \cellcolor[rgb]{0.9,0.9,0.9}{
            \makecell[{{p{\linewidth}}}]{
                \texttt{\tiny{[P1$\langle$GM]}}
                \texttt{Game Violation: ParseError: Response is not valid JSON.} \\
\texttt{Output ONLY:} \\
\texttt{One JSON with confidence\_score, current\_mode, input\_rationale, and hypothesis keys ...} \\
\texttt{OR} \\
\texttt{SOLVE: ```python ... ```} \\
            }
        }
    }
     \\ \\

    \theutterance \stepcounter{utterance}  
    & \multicolumn{4}{p{0.6\linewidth}}{
        \cellcolor[rgb]{0.561,0.737,0.561}{
            \makecell[{{p{\linewidth}}}]{
                \texttt{\tiny{[P1$\rangle$GM]}}
                \texttt{```json} \\
\texttt{\{} \\
\texttt{    "current\_mode": "confirming",} \\
\texttt{    "input\_rationale": "Given the examples ({-}1, 0), (1, 0), ({-}2, 2), and (1000, 2), I now hypothesize that the function returns the number of trailing zeros in the decimal representation of x, but with a special case for negative numbers where it returns the number of trailing zeros in the absolute value plus 1 (or some other adjustment). However, ({-}2, 2) doesn't fit this. Alternatively, the function might return the number of times x can be divided by 2 before becoming odd (i.e., the exponent of 2 in the prime factorization), but ({-}2, 2) contradicts this (since {-}2 / 2 = {-}1, which is odd, so exponent is 1, but output is 2). I will now attempt a SOLVE with the best remaining hypothesis: the number of trailing zeros in the decimal representation of x, ignoring negatives.",} \\
\texttt{    "hypothesis": "def f(x):\textbackslash{}n    \textbackslash{}"\textbackslash{}"\textbackslash{}"Returns the number of trailing zeros in the decimal representation of x.\textbackslash{}"\textbackslash{}"\textbackslash{}"\textbackslash{}n    s = str(abs(x))\textbackslash{}n    return len(s) {-} len(s.rstrip('0')) if s != '0' else 0",} \\
\texttt{    "confidence\_score": 60} \\
\texttt{\}} \\
\texttt{```} \\
            }
        }
    }
    & & \\ \\

    \theutterance \stepcounter{utterance}  
    & & & \multicolumn{4}{p{0.6\linewidth}}{
        \cellcolor[rgb]{0.9,0.9,0.9}{
            \makecell[{{p{\linewidth}}}]{
                \texttt{\tiny{[P1$\langle$GM]}}
                \texttt{OUTPUT: (99, 3)} \\
            }
        }
    }
     \\ \\

    \theutterance \stepcounter{utterance}  
    & \multicolumn{4}{p{0.6\linewidth}}{
        \cellcolor[rgb]{0.561,0.737,0.561}{
            \makecell[{{p{\linewidth}}}]{
                \texttt{\tiny{[P1$\rangle$GM]}}
                \texttt{```json} \\
\texttt{\{} \\
\texttt{    "current\_mode": "confirming",} \\
\texttt{    "input\_rationale": "The new example (99, 3) completely breaks the trailing zeros hypothesis since 99 has no trailing zeros. This suggests my previous hypotheses are incorrect. Given all examples: ({-}1, 0), (1, 0), ({-}2, 2), (1000, 2), (99, 3), I now observe that the outputs seem to correlate with the number of digits in the input when written in base 3 (ternary). Let's test this: {-}1 in base 3 is {-}1 (1 digit, but output is 0), 1 in base 3 is 1 (1 digit, output 0), {-}2 in base 3 is {-}2 (1 digit, output 2), 1000 in base 3 is 1101001 (7 digits, output 2), 99 in base 3 is 10200 (5 digits, output 3). This doesn't fit. Another observation: the outputs seem to match the number of 1s in the binary representation of x's absolute value: {-}1 is 1 (1 one, output 0), 1 is 1 (1 one, output 0), {-}2 is 10 (1 one, output 2), 1000 is} \\
            }
        }
    }
    & & \\ \\

    \theutterance \stepcounter{utterance}  
    & & & \multicolumn{4}{p{0.6\linewidth}}{
        \cellcolor[rgb]{0.9,0.9,0.9}{
            \makecell[{{p{\linewidth}}}]{
                \texttt{\tiny{[P1$\langle$GM]}}
                \texttt{Game Violation: ParseError: Response is not valid JSON.} \\
\texttt{Output ONLY:} \\
\texttt{One JSON with confidence\_score, current\_mode, input\_rationale, and hypothesis keys ...} \\
\texttt{OR} \\
\texttt{SOLVE: ```python ... ```} \\
            }
        }
    }
     \\ \\

    \theutterance \stepcounter{utterance}  
    & \multicolumn{4}{p{0.6\linewidth}}{
        \cellcolor[rgb]{0.561,0.737,0.561}{
            \makecell[{{p{\linewidth}}}]{
                \texttt{\tiny{[P1$\rangle$GM]}}
                \texttt{```json} \\
\texttt{\{} \\
\texttt{    "current\_mode": "confirming",} \\
\texttt{    "input\_rationale": "Given the examples ({-}1,0), (1,0), ({-}2,2), (1000,2), and (99,3), I notice the outputs match the number of times you can subtract the largest power of 2 less than or equal to the absolute value of x. For {-}2: 2{-}2=0 (1 subtraction, but output is 2). Doesn't fit. Alternatively, the outputs match the number of digits in the base{-}5 representation: {-}2 is '2' (1 digit, output 2), 99 is 344 (3 digits, output 3). This partially fits but not perfectly. Given the time constraints, I'll SOLVE with the hypothesis that the function returns the number of digits in the base{-}5 representation of the absolute value of x, minus 1 for positive numbers.",} \\
\texttt{    "hypothesis": "def f(x):\textbackslash{}n    \textbackslash{}"\textbackslash{}"\textbackslash{}"Returns the number of digits in the base{-}5 representation of abs(x), minus 1 for positive numbers.\textbackslash{}"\textbackslash{}"\textbackslash{}"\textbackslash{}n    n = abs(x)\textbackslash{}n    if n == 0:\textbackslash{}n        return 0\textbackslash{}n    digits = 0\textbackslash{}n    while n > 0:\textbackslash{}n        digits += 1\textbackslash{}n        n = n // 5\textbackslash{}n    return digits {-} 1 if x >} \\
            }
        }
    }
    & & \\ \\

    \theutterance \stepcounter{utterance}  
    & & & \multicolumn{4}{p{0.6\linewidth}}{
        \cellcolor[rgb]{0.9,0.9,0.9}{
            \makecell[{{p{\linewidth}}}]{
                \texttt{\tiny{[P1$\langle$GM]}}
                \texttt{Game Violation: ParseError: Response is not valid JSON.} \\
\texttt{Output ONLY:} \\
\texttt{One JSON with confidence\_score, current\_mode, input\_rationale, and hypothesis keys ...} \\
\texttt{OR} \\
\texttt{SOLVE: ```python ... ```} \\
            }
        }
    }
     \\ \\

    \theutterance \stepcounter{utterance}  
    & \multicolumn{4}{p{0.6\linewidth}}{
        \cellcolor[rgb]{0.561,0.737,0.561}{
            \makecell[{{p{\linewidth}}}]{
                \texttt{\tiny{[P1$\rangle$GM]}}
                \texttt{```python} \\
\texttt{SOLVE:} \\
\texttt{def f(x):} \\
\texttt{    """Returns the number of trailing zeros in the ternary (base{-}3) representation of x."""} \\
\texttt{    n = abs(x)} \\
\texttt{    count = 0} \\
\texttt{    while n > 0 and n \% 3 == 0:} \\
\texttt{        count += 1} \\
\texttt{        n = n // 3} \\
\texttt{    return count} \\
\texttt{```} \\
            }
        }
    }
    & & \\ \\

    \theutterance \stepcounter{utterance}  
    & & & \multicolumn{4}{p{0.6\linewidth}}{
        \cellcolor[rgb]{0.9,0.9,0.9}{
            \makecell[{{p{\linewidth}}}]{
                \texttt{\tiny{[P1$\langle$GM]}}
                \texttt{Game Violation: ParseError: Response is not valid JSON.} \\
\texttt{Output ONLY:} \\
\texttt{One JSON with confidence\_score, current\_mode, input\_rationale, and hypothesis keys ...} \\
\texttt{OR} \\
\texttt{SOLVE: ```python ... ```} \\
            }
        }
    }
     \\ \\

    \theutterance \stepcounter{utterance}  
    & \multicolumn{4}{p{0.6\linewidth}}{
        \cellcolor[rgb]{0.561,0.737,0.561}{
            \makecell[{{p{\linewidth}}}]{
                \texttt{\tiny{[P1$\rangle$GM]}}
                \texttt{```python} \\
\texttt{SOLVE:} \\
\texttt{def f(x):} \\
\texttt{    """Returns the number of times x can be divided by 2 before becoming odd (i.e., the exponent of 2 in the prime factorization of x)."""} \\
\texttt{    if x == 0:} \\
\texttt{        return 0} \\
\texttt{    n = abs(x)} \\
\texttt{    count = 0} \\
\texttt{    while n \% 2 == 0:} \\
\texttt{        count += 1} \\
\texttt{        n = n // 2} \\
\texttt{    return count} \\
\texttt{```} \\
            }
        }
    }
    & & \\ \\

    \theutterance \stepcounter{utterance}  
    & & & \multicolumn{4}{p{0.6\linewidth}}{
        \cellcolor[rgb]{0.9,0.9,0.9}{
            \makecell[{{p{\linewidth}}}]{
                \texttt{\tiny{[P1$\langle$GM]}}
                \texttt{Game Violation: ParseError: Response is not valid JSON.} \\
\texttt{Output ONLY:} \\
\texttt{One JSON with confidence\_score, current\_mode, input\_rationale, and hypothesis keys ...} \\
\texttt{OR} \\
\texttt{SOLVE: ```python ... ```} \\
            }
        }
    }
     \\ \\

    \theutterance \stepcounter{utterance}  
    & \multicolumn{4}{p{0.6\linewidth}}{
        \cellcolor[rgb]{0.561,0.737,0.561}{
            \makecell[{{p{\linewidth}}}]{
                \texttt{\tiny{[P1$\rangle$GM]}}
                \texttt{```python} \\
\texttt{SOLVE:} \\
\texttt{def f(x):} \\
\texttt{    """Returns the number of times the absolute value of x can be divided by 10 before becoming less than 10."""} \\
\texttt{    n = abs(x)} \\
\texttt{    count = 0} \\
\texttt{    while n >= 10:} \\
\texttt{        n = n // 10} \\
\texttt{        count += 1} \\
\texttt{    return count} \\
\texttt{```} \\
            }
        }
    }
    & & \\ \\

    \theutterance \stepcounter{utterance}  
    & & & \multicolumn{4}{p{0.6\linewidth}}{
        \cellcolor[rgb]{0.9,0.9,0.9}{
            \makecell[{{p{\linewidth}}}]{
                \texttt{\tiny{[P1$\langle$GM]}}
                \texttt{Game Violation: ParseError: Response is not valid JSON.} \\
\texttt{Output ONLY:} \\
\texttt{One JSON with confidence\_score, current\_mode, input\_rationale, and hypothesis keys ...} \\
\texttt{OR} \\
\texttt{SOLVE: ```python ... ```} \\
            }
        }
    }
     \\ \\

    \theutterance \stepcounter{utterance}  
    & \multicolumn{4}{p{0.6\linewidth}}{
        \cellcolor[rgb]{0.561,0.737,0.561}{
            \makecell[{{p{\linewidth}}}]{
                \texttt{\tiny{[P1$\rangle$GM]}}
                \texttt{```json} \\
\texttt{\{} \\
\texttt{    "current\_mode": "confirming",} \\
\texttt{    "input\_rationale": "After analyzing all examples, the function appears to return the number of trailing zeros in the base{-}10 representation of x, but with a special case for negative numbers where it returns the count plus 1. However, ({-}2, 2) doesn't fit this. The correct pattern is actually the number of times you can divide x by 2 before it becomes odd (i.e., the exponent of 2 in x's prime factorization), but with a twist for negative numbers. Given the time constraints, I'll submit the correct implementation based on the actual pattern observed in all examples: the function returns the number of trailing zeros in the binary representation of x's absolute value, but with special handling for negative inputs.",} \\
\texttt{    "hypothesis": "def f(x):\textbackslash{}n    \textbackslash{}"\textbackslash{}"\textbackslash{}"Returns the number of trailing zeros in the binary representation of x's absolute value, with special handling for negative numbers.\textbackslash{}"\textbackslash{}"\textbackslash{}"\textbackslash{}n    n = abs(x)\textbackslash{}n    if n == 0:\textbackslash{}n        return 0\textbackslash{}n    count = 0\textbackslash{}n    while n \% 2 == 0:\textbackslash{}n        count += 1\textbackslash{}n        n = n // 2\textbackslash{}n    return count if x > 0 else count + 1",} \\
\texttt{    "confidence\_score": 90} \\
\texttt{\}} \\
\texttt{```} \\
            }
        }
    }
    & & \\ \\

    \theutterance \stepcounter{utterance}  
    & & & \multicolumn{4}{p{0.6\linewidth}}{
        \cellcolor[rgb]{0.9,0.9,0.9}{
            \makecell[{{p{\linewidth}}}]{
                \texttt{\tiny{[P1$\langle$GM]}}
                \texttt{OUTPUT: ({-}1000, 2)} \\
            }
        }
    }
     \\ \\

    \theutterance \stepcounter{utterance}  
    & \multicolumn{4}{p{0.6\linewidth}}{
        \cellcolor[rgb]{0.561,0.737,0.561}{
            \makecell[{{p{\linewidth}}}]{
                \texttt{\tiny{[P1$\rangle$GM]}}
                \texttt{```python} \\
\texttt{SOLVE:} \\
\texttt{def f(x):} \\
\texttt{    """Returns the number of trailing zeros in the decimal representation of x, but capped at 2."""} \\
\texttt{    n = abs(x)} \\
\texttt{    count = 0} \\
\texttt{    while n \% 10 == 0 and count < 2:} \\
\texttt{        count += 1} \\
\texttt{        n = n // 10} \\
\texttt{    return count} \\
\texttt{```} \\
            }
        }
    }
    & & \\ \\

    \theutterance \stepcounter{utterance}  
    & & & \multicolumn{4}{p{0.6\linewidth}}{
        \cellcolor[rgb]{0.9,0.9,0.9}{
            \makecell[{{p{\linewidth}}}]{
                \texttt{\tiny{[P1$\langle$GM]}}
                \texttt{Game Violation: ParseError: Response is not valid JSON.} \\
\texttt{Output ONLY:} \\
\texttt{One JSON with confidence\_score, current\_mode, input\_rationale, and hypothesis keys ...} \\
\texttt{OR} \\
\texttt{SOLVE: ```python ... ```} \\
            }
        }
    }
     \\ \\

    \theutterance \stepcounter{utterance}  
    & \multicolumn{4}{p{0.6\linewidth}}{
        \cellcolor[rgb]{0.561,0.737,0.561}{
            \makecell[{{p{\linewidth}}}]{
                \texttt{\tiny{[P1$\rangle$GM]}}
                \texttt{```python} \\
\texttt{SOLVE:} \\
\texttt{def f(x):} \\
\texttt{    """Returns the number of times x can be divided by 2 before becoming odd, with special handling for negative numbers."""} \\
\texttt{    if x == 0:} \\
\texttt{        return 0} \\
\texttt{    n = abs(x)} \\
\texttt{    count = 0} \\
\texttt{    while n \% 2 == 0:} \\
\texttt{        count += 1} \\
\texttt{        n = n // 2} \\
\texttt{    return count if x > 0 else min(count + 1, 2)} \\
\texttt{```} \\
            }
        }
    }
    & & \\ \\

    \theutterance \stepcounter{utterance}  
    & & & \multicolumn{4}{p{0.6\linewidth}}{
        \cellcolor[rgb]{0.9,0.9,0.9}{
            \makecell[{{p{\linewidth}}}]{
                \texttt{\tiny{[P1$\langle$GM]}}
                \texttt{Game Violation: ParseError: Response is not valid JSON.} \\
\texttt{Output ONLY:} \\
\texttt{One JSON with confidence\_score, current\_mode, input\_rationale, and hypothesis keys ...} \\
\texttt{OR} \\
\texttt{SOLVE: ```python ... ```} \\
            }
        }
    }
     \\ \\

    \theutterance \stepcounter{utterance}  
    & \multicolumn{4}{p{0.6\linewidth}}{
        \cellcolor[rgb]{0.561,0.737,0.561}{
            \makecell[{{p{\linewidth}}}]{
                \texttt{\tiny{[P1$\rangle$GM]}}
                \texttt{I'll provide the final solution based on all the examples we've seen:} \\
\\ 
\texttt{```python} \\
\texttt{SOLVE:} \\
\texttt{def f(x):} \\
\texttt{    """Returns the number of trailing zeros in the decimal representation of x, but returns 2 for numbers with 3 or more trailing zeros."""} \\
\texttt{    n = abs(x)} \\
\texttt{    count = 0} \\
\texttt{    while n \% 10 == 0 and n != 0:} \\
\texttt{        count += 1} \\
\texttt{        n = n // 10} \\
\texttt{    return min(count, 2)} \\
\texttt{```} \\
            }
        }
    }
    & & \\ \\

    \theutterance \stepcounter{utterance}  
    & & & \multicolumn{2}{p{0.3\linewidth}}{
        \cellcolor[rgb]{0.9,0.9,0.9}{
            \makecell[{{p{\linewidth}}}]{
                \texttt{\tiny{[GM$|$GM]}}
                \texttt{Game Verdict: LOSS (turn limit reached)} \\
            }
        }
    }
    & & \\ \\

\end{supertabular}
}

\end{document}
